\section{Trie / Booth}
\subsection{Trie (Binary)}
\begin{minted}{c++}
// Initialize: Trie T;
// Insert/Delete: T.add(val, 1) / T.add(val, -1);
struct Trie {
  static const int B = 31;
  vector<array<int, 2>> nxt;
  vector<int> cnt;
  Trie() { nxt.push_back({0, 0}); cnt.push_back(0); }
  void add(int x, int v = 1) {
    int u = 0;
    cnt[u] += v;
    for (int i = B - 1; i >= 0; --i) {
      int d = (x >> i) & 1;
      if (!nxt[u][d]) {
        nxt[u][d] = nxt.size();
        nxt.push_back({0, 0});
        cnt.push_back(0);
      }
      u = nxt[u][d];
      cnt[u] += v;
    }
  }
  int count(int x) {
    int u = 0;
    for (int i = B - 1; i >= 0; --i) {
      int d = (x >> i) & 1;
      if (!nxt[u][d] || !cnt[nxt[u][d]]) return 0;
      u = nxt[u][d];
    }
    return cnt[u];
  }
  int count_xor_less(int x, int k) {
    int u = 0, ans = 0;
    for (int i = B - 1; i >= 0 && u; --i) {
      int d = (x >> i) & 1;
      if ((k >> i) & 1) {
        if (nxt[u][d]) ans += cnt[nxt[u][d]];
        u = nxt[u][d ^ 1];
      } else u = nxt[u][d];
    }
    return ans;
  }
  int get_kth(int k) {
    if (cnt[0] < k) return -1;
    int u = 0, ans = 0;
    for (int i = B - 1; i >= 0; --i) {
      int left = nxt[u][0] ? cnt[nxt[u][0]] : 0;
      if (k <= left) u = nxt[u][0];
      else k -= left, u = nxt[u][1], ans |= (1 << i);
    }
    return ans;
  }
  int get_xor_kth(int x, int k) {
    if (cnt[0] < k) return -1;
    int u = 0, ans = 0;
    for (int i = B - 1; i >= 0; --i) {
      int d = (x >> i) & 1;
      int req = nxt[u][d] ? cnt[nxt[u][d]] : 0;
      if (k <= req) u = nxt[u][d], ans |= (d << i);
      else k -= req, u = nxt[u][d ^ 1], ans |= ((d ^ 1) << i);
    }
    return ans;
  }
  int64_t go(int u, int b, int x, int t1, int t2) {
    if (!u || !t1 || t2 == 2) return 0;
    if ((t1 == 2 && !t2) || b < 0) return cnt[u];
    int d = (x >> b) & 1;
    int64_t res = 0;
    if (d) {
      res += go(nxt[u][1], b - 1, x, t1 == 1 ? 0 : t1, t2 == 1 ? 0 : t2);
      res += go(nxt[u][0], b - 1, x, t1 == 1 ? 2 : t1, t2);
    } else {
      res += go(nxt[u][1], b - 1, x, t1, t2 == 1 ? 2 : t2);
      res += go(nxt[u][0], b - 1, x, t1, t2);
    }
    return res;
  }
  int64_t count_special(int x) { return go(0, B - 1, x, 1, 1); }
};
\end{minted}
\subsection{Booth's Algorithm (Min Lex Rotation)}
\begin{minted}{c++}
int minLex(string &s) {
  int len = s.size();
  int n = 2 * len, i = 0, j = 1, k = 0;
  int a, b;
  while (i + k < n && j + k < n) {
    a = (i + k >= len) ? s[i + k - len] : s[i + k];
    b = (j + k >= len) ? s[j + k - len] : s[j + k];
    if (a == b) ++k;
    else if (a > b) {
      i = i + k + 1;
      if (i <= j) i = j + 1;
      k = 0;
    } else {
      j = j + k + 1;
      if (j <= i) j = i + 1;
      k = 0;
    }
  }
  return min(i, j);
}
\end{minted}
\subsection{GCD Properties}
In particular, recalling that gcd is a positive integer valued function we obtain that $\gcd(a, b \cdot c) = 1$ if and only if $\gcd(a, b) = 1$ and $\gcd(a, c) = 1$.

15. $\gcd(a_1, \ldots, a_r) = \gcd(\gcd(a_1, \ldots, a_{r-1}), a_r)$ when none $=0$

16. $\gcd(a,b) \cdot \text{lcm}(a,b) = |a \cdot b|$, so for any two numbers LCM is always multiple of GCD.

17. In a Cartesian coordinate system, $\gcd(a, b)$ can be interpreted as the number of segments between points with integral coordinates on the straight line segment joining the points $(0, 0)$ and $(a, b)$.

18. Two numbers are called relatively prime, or coprime, if their greatest common divisor equals 1. For example, 9 and 28 are relatively prime.

19. Any two consecutive integers will be Co-Prime. Ex: 25,26 is co-prime. 5,6 is co-prime and so on. This will work because, let x divides 25 so remainder is something. If we divided 26 by x then remainder will be 1. This is the logic.

20. Two or more odd consecutive positive integer GCD is also 1.
\begin{minted}{c++}
int gcdFast(int a, int b)
{
  int result;
  __asm__ __volatile__(
    "movl %1, %%eax;"
    "movl %2, %%ebx;"
    "CONTD: cmpl $0, %%ebx;"
    "je DONE;"
    "xorl %%edx, %%edx;"
    "idivl %%ebx;"
    "movl %%ebx, %%eax;"
    "movl %%edx, %%ebx;"
    "jmp CONTD;"
    "DONE: movl %%eax, %0;"
    : "=g"(result)
    : "g"(a), "g"(b));
  return result;
}
\end{minted}
\subsection{Extended Euclid and CRT}
\begin{minted}{c++}
using T = __int128_t;
T egcd(T a, T b, T &x, T &y) {
  T xx = y = 0, yy = x = 1;
  while (b) {
    T q = a / b;
    T t = b; b = a % b; a = t;
    t = xx; xx = x - q * xx; x = t;
    t = yy; yy = y - q * yy; y = t;
  }
  return a;
}
// finds x such that x % m1 = a1, x % m2 = a2.
// m1 and m2 may not be coprime
// here, x is unique modulo m = lcm(m1, m2).
// returns (x, m). on failure, m = -1.
pair<T, T> crt(T a1, T m1, T a2, T m2) {
  T p, q;
  T g = egcd(m1, m2, p, q);
  if (a1 % g != a2 % g) return {0, -1};
  T m = m1 / g * m2;
  p = (p % m + m) % m;
  q = (q % m + m) % m;
  T x = p * a2 % m * (m1 / g) % m + q * a1 % m *
  (m2 / g) % m;
  return {x % m, m};
}
\end{minted}
\subsection{SPF, Mobius, Primes in O(N)}
\begin{minted}{c++}
const int N = 1e6 + 2;
vector<int> spf(N), mob(N), primes;
void build() {
  mob[1] = 1;
  for (int i = 2; i < N; i++) {
    if (spf[i] == 0) {
      spf[i] = i, mob[i] = -1,
      primes.push_back(i);
    }
    for (int j = 0; i * primes[j] < N; j++) {
      spf[i * primes[j]] = primes[j];
      if (i % primes[j])
        mob[i * primes[j]] = mob[i] *
        mob[primes[j]];
      else mob[i * primes[j]] = 0;
      if (primes[j] == spf[i]) break;
    }
  }
}
\end{minted}
\subsection{Segmented Sieve}
\begin{minted}{c++}
vector<char> segmentedSieve(long long L, long
long R) {
  // generate all primes up to sqrt(R)
  long long lim = sqrt(R);
  vector<char> mark(lim + 1, false);
  vector<long long> primes;
  for (long long i = 2; i <= lim; ++i) {
    if (!mark[i]) {
      primes.emplace_back(i);
      for (long long j = i * i; j <= lim; j += i)
        mark[j] = true;
    }
  }
  vector<char> isPrime(R - L + 1, true);
  for (long long i : primes) {
    ll sm = (L / i) * i;
    if (sm < L) sm += i;
    for (long long j = max(i * i, sm); j <= R; j
    += i)
      isPrime[j - L] = false;
  }
  if (L == 1) isPrime[0] = false;
  return isPrime;
}
\end{minted}
\subsection{Phi (Euler Totient)}
\begin{minted}{c++}
int phi(int n) {
  int result = n;
  for (int i = 2; i * i <= n; i++) {
    if (n % i == 0) {
      while (n % i == 0)
        n /= i;
      result -= result / i;
    }
  }
  if (n > 1)
    result -= result / n;
  return result;
}
void phi_1_to_n(){
  phi[1] = 1;
  for(int i = 2; i < N; i++) {
    if(phi[i] == 0){
      primes.push_back(i);
      phi[i] = i - 1;
    }
    for(int p : primes){
      if(i * p >= N) break;
      if(i % p == 0){
        phi[i * p] = phi[i] * p;
        break;
      }else {
        phi[i * p] = phi[i] * (p - 1);
      }
    }
  }
}
\end{minted}
\subsection{Miller-Rabin Primality Test}
\begin{minted}{c++}
ull pow(ull a, ull t, ull mod) {
  ull r = 1;
  for (; t; t >>= 1, a = (ull)a * a % mod)
    if (t & 1) r = (ull)r * a % mod;
  return r;
}
bool isprime(ull n) {
  if (n < 2 || n % 6 % 4 != 1) return (n | 1) == 3;
  ull A[] = {2, 325, 9375, 28178, 450775,
  9780504, 1795265022},
  s = __builtin_ctzll(n - 1), d = n >> s;
  for (ull a : A) {
    ull p = pow(a % n, d, n), i = s;
    while (p != 1 && p != n - 1 && a % n && i--)
      p = (ull)p * p % n;
    if (p != n - 1 && i != s) return 0;
  }
  return 1;
}
\end{minted}
\subsection{Matrix Exponentiation}
\begin{minted}{c++}
vector<vector<int>>
multiply(vector<vector<int>>& a,
vector<vector<int>>& b) {
  int n = sz(a), m = sz(a[0]), p = sz(b[0]);
  vector<vector<int>> result(n, vector<int>(p,
  0));
  for (int i = 0; i < n; i++) {
    for (int k = 0; k < m; k++) {
      if (a[i][k] == 0) continue;
      for (int j = 0; j < p; j++) {
        result[i][j] = ((result[i][j] + a[i][k] *
        b[k][j]) % mod + mod) % mod;
      }
    }
  }
  return result;
}
vector<vector<int>> expo(vector<vector<int>> t,
int p) {
  int n = t.size();
  vector<vector<int>> result(n, vector<int>(n,
  0));
  for (int i = 0; i < n; i++) {
    result[i][i] = 1;
  }
  while (p > 0) {
    if (p & 1) result = multiply(result, t);
    t = multiply(t, t);
    p >>= 1;
  }
  return result;
}
\end{minted}
